% !TeX root=../main.tex
\chapter{مدل سیستم و صورت‌بندی ریاضی مسئله}
\label{chap:system-model}

\section{مقدمه}
فصل قبل مبانی و پیشینه را مرور کرد؛ در این فصل مسئلهٔ تحقیق به‌صورت دقیق مدل‌سازی و پارامترهای مجهول تعریف می‌شوند.

\section{ساختار عمومی کد توربو}
ساختار موازی متداول شامل دو کدگذار RSC، درهم‌ریز داخلی، و در صورت وجود درهم‌ریز خارجی و الگوی پانچینگ توصیف می‌شود.

\section{مدل کدگذار RSC}
\subsection{مدل چندجمله‌ای}
\subsection{تابع انتقال}
\subsection{نمایش باینری و اکتال}
\subsection{معادلات بازگشتی}
روابط بازگشتی میان بیت سیستماتیک و پاریتی، قیود جبری لازم برای بازسازی را فراهم می‌کنند.

\section{مدل درهم‌ریز داخلی}
درهم‌ریز داخلی به‌صورت جایگزشت روی بلوک اطلاعات مدل می‌شود و فضای حالت آن در حالت کلی بسیار بزرگ است.

\section{مدل درهم‌ریز خارجی}
درهم‌ریز خارجی چینش بیت‌های خروجی را تعیین می‌کند و باید در مجموعهٔ پارامترهای مجهول لحاظ شود.

\section{مدل پانچینگ}
الگوی پانچینگ نرخ مؤثر کد را تغییر می‌دهد؛ در مدل پایه ممکن است معلوم فرض شود و در فصل توسعه به‌صورت ناشناخته بررسی گردد.

\section{مدل خاتمه ترلیس}
با درنظرگرفتن بیت‌های خاتمه، طول فریم پیش از درهم‌ریز خارجی می‌تواند از رابطه‌ای مانند
\begin{equation}
	N = 3L + 4m
	\label{eq:frame-length}
\end{equation}
به‌دست آید که در آن \(L\) طول بلوک اطلاعات و \(m\) حافظهٔ کدگذار است.

\section{مدل کانال}
\subsection{AWGN}
\subsection{BPSK}
\subsection{BSC معادل}
پس از آشکارسازی سخت، مشاهده به‌صورت دنباله‌های باینری نویزی در مدل BSC توصیف می‌شود.

\section{مدل دریافت}
فرض‌های مشاهده شامل مشخص‌بودن یا نبودن مرز فریم، نوع آشکارسازی سخت/نرم، تعداد فریم‌های دریافتی \(F\)، نرخ خطای کانال، و ناشناخته‌بودن ساختار کد است.

\section{تعریف پارامترهای مجهول}
مجموعهٔ پارامترهای مجهول به‌صورت
\begin{equation}
	\theta = \{m, g_0(D), g_1(D), \pi_i, \pi_o, L\}
	\label{eq:theta}
\end{equation}
تعریف می‌شود که شامل حافظه، پلی‌نوم‌های مولد، درهم‌ریزهای داخلی و خارجی، و طول بلوک است.

\section{تعریف مسئله بازسازی کور}
با دریافت مشاهدات \(R\)، هدف تخمین مجموعهٔ مقادیر مجاز \(\theta\) است. خروجی الگوریتم لزوماً یک پاسخ یکتا نیست و ممکن است مجموعه‌ای از بازسازی‌های ساختاریاً معادل (ابهام ساختاری) به‌دست آید.
